\documentclass[11pt,a4paper]{article}

\usepackage[T1]{fontenc}
\usepackage[utf8]{inputenc}
\usepackage{amsmath,amssymb,amsthm}
\usepackage{geometry}
\usepackage{xcolor}

\geometry{margin=1in}
\setlength{\parindent}{0pt}
\setlength{\parskip}{0.62em}
\linespread{1.035}

\newtheoremstyle{soft}{0.7em}{0.7em}{\itshape}{}{\bfseries\color{blue!35!black}}{.}{0.5em}{}
\theoremstyle{soft}
\newtheorem{proposition}{Proposition}[section]
\newtheorem{conjecture}[proposition]{Conjecture}
\theoremstyle{definition}
\newtheorem{definition}[proposition]{Definition}
\newtheorem{remark}[proposition]{Remark}

\newcommand{\C}{\mathbb C}
\newcommand{\Q}{\mathbf Q}
\newcommand{\partline}[2]{\section*{#1\quad #2}\addcontentsline{toc}{section}{#1: #2}}
\newcommand{\idea}[1]{\begin{center}\fcolorbox{blue!20!black}{blue!4}{\begin{minipage}{0.90\linewidth}\small #1\end{minipage}}\end{center}}
\newcommand{\status}[1]{\begin{center}\fcolorbox{black!25}{black!3}{\begin{minipage}{0.90\linewidth}\small\textbf{Status.} #1\end{minipage}}\end{center}}

\title{Universitas Pandrosion\\[-0.2em]\large Critical Quotients, Fibers, and a Program around Smale's Mean Value Conjecture}
\author{Ivan Besevic}
\date{Revised compact edition, May 2026}

\begin{document}
\maketitle

\begin{abstract}
This manuscript rewrites the nine Pandrosion--Smale notes as a compact research program.  The guiding object is the finite slope
\[
   Q_P(a,b)=\frac{P(b)-P(a)}{b-a},
\]
and, for polynomial systems, a matrix slope satisfying
\[
   F(b)-F(a)=\Q_F(a,b)(b-a).
\]
The finite-slope identity itself belongs to the classical theory of divided differences and slope matrices.  The Pandrosion viewpoint developed here is therefore not the claim that this identity is new.  Rather, the claim is geometric: critical points, fibers, scaling, starts, and line searches can be organized around the same exact quotient.  Algebraic identities are separated from programmatic or experimental statements.
\end{abstract}

\tableofcontents

\section*{Reading principle}
\addcontentsline{toc}{section}{Reading principle}
Pandrosion is a language of exact comparison.  Newton differentiates at one point.  Pandrosion compares two points and asks for the slope that makes the comparison exact.  In one variable this is a quotient; in several variables it is a finite-slope matrix.  This note uses that quotient as a geometric vocabulary for Smale's Mean Value Conjecture and for derivative-free root finding.

\idea{The useful separation is this: the slope identity is classical; the Pandrosion program is the geometric organization built around it.  Parts I--VI concern Smale's critical-point inequality.  Parts VII--IX describe algorithmic ideas whose rigorous finite-slope convergence and conditioned complexity are treated in companion work.}

\begin{figure}[h]
\centering
\setlength{\unitlength}{0.75mm}
\begin{picture}(95,43)
  \put(5,6){\vector(1,0){82}}
  \put(8,4){$z$}
  \put(10,4){\vector(0,1){35}}
  \put(12,38){$P(z)$}
  \qbezier(18,12)(38,34)(70,20)
  \qbezier(70,20)(78,18)(85,32)
  \put(25,15){\circle*{2.4}}\put(27,12){$a$}
  \put(75,27){\circle*{2.4}}\put(68,31){$b$}
  \linethickness{0.45mm}\put(25,15){\line(5,3){50}}
  \put(45,29){$Q_P(a,b)$}
  \linethickness{0.2mm}\multiput(25,6)(0,3){4}{\line(0,1){1.5}}
  \multiput(75,6)(0,3){8}{\line(0,1){1.5}}
\end{picture}
\caption{The basic Pandrosion quotient is the exact slope between two finite points.  Newton appears only in the limit $b\to a$.}
\end{figure}

\partline{Part I}{The finite-slope operator}

Let $P\in\C[z]$.  For $a\ne b$, define
\begin{equation}\label{eq:scalar-Q}
   Q_P(a,b)=\frac{P(b)-P(a)}{b-a}.
\end{equation}
If $b\to a$, then $Q_P(a,b)\to P'(a)$.  Thus the derivative is an infinitesimal limit of an exact finite comparison.

For systems $F:\C^n\to\C^n$, one uses a slope matrix.
\begin{definition}[Finite-slope matrix]
A matrix $\Q_F(a,b)\in\C^{n\times n}$ is a finite-slope matrix between $a,b\in\C^n$ if
\begin{equation}\label{eq:matrix-Q}
   F(b)-F(a)=\Q_F(a,b)(b-a).
\end{equation}
For polynomial systems a canonical coordinate-path choice is obtained by expanding each monomial telescopically.
\end{definition}

\begin{remark}[Classical source]
Matrices satisfying \eqref{eq:matrix-Q} are part of the classical theory of divided differences, secant methods, and slope matrices; see Ortega--Rheinboldt \cite{ortegarheinboldt} and related slope/secant literature \cite{potraptak,krawczyk,neumaier}.  The Pandrosion contribution in this series is the geometric use of these finite slopes: the choice of probes, charts, scaling, fibers, and starts.
\end{remark}

Once $\Q_F(a,b)$ is available, the fixed-anchor Pandrosion step is
\begin{equation}\label{eq:pandrosion-step}
   \mathcal P_{F,a}(b)=a-\Q_F(a,b)^{-1}F(a),
\end{equation}
whenever the matrix is invertible.  In the auto-probe limit $b=a$, this becomes Newton.  Away from that limit it is a derivative-free finite-slope correction.

The original one-dimensional root-extraction formula appears as a special case.  Let
\[
   S_p(s)=1+s+\cdots+s^{p-1}.
\]
For $x\ne0$ and $p\ge2$, the map
\begin{equation}\label{eq:root-map}
   h_{p,x}(s)=1-\frac{x-1}{xS_p(s)}
\end{equation}
has fixed points $s=x^{-1/p}$, one for each branch.  The finite sum $S_p$ is the monomial finite slope in disguise:
\[
   b^p-a^p=(b-a)\sum_{k=0}^{p-1}b^{p-1-k}a^k.
\]
This is the scalar model from which the higher-dimensional construction is lifted.

\status{The local convergence and complexity of \eqref{eq:pandrosion-step} are not proved in this program note.  They are treated in the companion finite-slope theorem paper.  Here the operator is introduced as the common language for the Smale and fiber identities below.}

\partline{Part II}{Smale's Mean Value Conjecture as a quotient statement}

Smale's Mean Value Conjecture concerns a polynomial $P$ of degree $d$, a point $z$ with $P'(z)\ne0$, and the critical points $c$ of $P$.  In finite-slope notation the conjectural inequality becomes
\begin{equation}\label{eq:smale-mvc}
   \min_{P'(c)=0}\left|\frac{Q_P(z,c)}{P'(z)}\right|
   \le \frac{d-1}{d}.
\end{equation}
Since $P'(z)=Q_P(z,z)$, this is a clean quotient translation: among the critical points, at least one secant quotient from $z$ is no larger than $(d-1)/d$ of the infinitesimal quotient at $z$.

\status{Equation \eqref{eq:smale-mvc} is a reformulation, not a proof and not a new equivalent difficulty reduction.  Its value is geometric: it turns the conjecture into a comparison between finite and infinitesimal slopes.}

For quadratics the statement is exact.  If $P(z)=(z-\alpha)(z-\beta)$ and $c=(\alpha+\beta)/2$, then
\[
   \left|\frac{Q_P(z,c)}{P'(z)}\right|=\frac12.
\]
For higher degree, the product over all critical points is controlled by a resultant.  If $c_1,\ldots,c_{d-1}$ are the roots of $P'$ and $R(w)=Q_P(z,w)$, then
\begin{equation}\label{eq:critical-product}
   \prod_{j=1}^{d-1}Q_P(z,c_j)=\frac1d\operatorname{Res}(R,P').
\end{equation}
The identity is elementary but useful: it packages all critical targets into one algebraic object.

\partline{Part III}{The vanishing identity}

Normalize at the origin:
\[
   P(0)=0,\qquad P'(0)=1.
\]
Write the roots as $0,\alpha_2,\ldots,\alpha_d$, and set
\[
   R(z)=\prod_{k=2}^{d}(z-\alpha_k),\qquad P(z)=zR(z).
\]
The partial-fraction expansion of $1/P(z)$ gives
\begin{equation}\label{eq:vanishing-identity}
   1=-\sum_{k=2}^{d}\frac{1}{\alpha_k R'(\alpha_k)}.
\end{equation}
This is a root-side constraint.  It does not solve the conjecture, but it explains why extremal configurations must be balanced across the nonzero roots.

At a critical point $c$, one has
\begin{equation}\label{eq:critical-R}
   R(c)=-cR'(c),
\end{equation}
since $P'(c)=R(c)+cR'(c)=0$.  Thus the Smale value can be read either from $R(c)$ or from the derivative of the root polynomial at $c$.

The symmetric family
\begin{equation}\label{eq:zd-z}
   P(z)=z^d-z
\end{equation}
attains the value $(d-1)/d$ exactly.  It is therefore the natural model extremal for the series.

\partline{Part IV}{Fibers and inverse Pandrosion quotients}

At a critical point $c$, the value $P(c)$ has a fiber.  Removing the critical point itself, call the remaining points the co-fiber of $c$.  The quotient from $c$ back to the origin is
\begin{equation}\label{eq:inverse-quotient}
   Q_P(c,0)=\frac{P(c)}{c}.
\end{equation}
Thus the Smale ratio is an inverse finite-slope quotient.

If the co-fiber points are $z_1,\ldots,z_{d-2}$, then
\begin{equation}\label{eq:cofiber-product}
   \left|\frac{P(c)}{c}\right|=|c|\prod_{j=1}^{d-2}|z_j|.
\end{equation}
A value estimate becomes a geometric estimate on a fiber.

The following two identities are also simple partial-fraction consequences.  Let $\alpha_0,\ldots,\alpha_{d-1}$ be the roots of $P$.  At a critical point $c$,
\begin{align}
   \sum_{j=0}^{d-1}\frac{1}{(\alpha_j-c)P'(\alpha_j)}&=-\frac{1}{P(c)},\label{eq:fiber-bridge}\\
   \sum_{j=0}^{d-1}\frac{1}{(\alpha_j-c)^2P'(\alpha_j)}&=0.\label{eq:fiber-vanish}
\end{align}
The first equation bridges the fiber to the value $P(c)$.  The second is an orthogonality relation.  Together they make precise the intuition that the conjecture is controlled by fiber geometry, not only by the locations of critical points.

\begin{figure}[h]
\centering
\setlength{\unitlength}{0.8mm}
\begin{picture}(90,50)
  \put(8,25){\vector(1,0){72}}
  \put(45,5){\vector(0,1){40}}
  \put(54,30){\circle*{2.4}}\put(57,31){$c$}
  \put(16,35){\circle*{2}}\put(12,38){$\alpha_1$}
  \put(26,14){\circle*{2}}\put(20,10){$\alpha_2$}
  \put(70,16){\circle*{2}}\put(71,12){$\alpha_3$}
  \put(76,37){\circle*{2}}\put(73,40){$\alpha_4$}
  \put(20,23){\circle*{2}}\put(42,38){\circle*{2}}\put(65,27){\circle*{2}}
  \linethickness{0.45mm}\put(54,30){\vector(-2,-1){32}}
  \put(28,17){$Q_P(c,0)$}
  \put(18,43){co-fiber}
\end{picture}
\caption{The inverse quotient $P(c)/c$ can be read through the fiber of the critical value.}
\end{figure}

For normalized polynomials
\[
   P(z)=z^d+a_{d-1}z^{d-1}+\cdots+a_1z,
\]
the critical equation gives the Pandrosion reduction
\begin{equation}\label{eq:pandrosion-reduction}
   \frac{P(c)}{c}=\widehat q(c),\qquad
   \widehat q(z)=\frac1d\widetilde q(z),\qquad
   \widetilde q(z)=\sum_{k=1}^{d-1}k a_{d-k}z^{d-1-k}.
\end{equation}
Thus Smale's bound becomes the search for a critical point satisfying $|\widehat q(c)|\le(d-1)/d$, equivalently $|\widetilde q(c)|\le d-1$.

\partline{Part V}{The quartic slice}

For degree four, write
\begin{equation}\label{eq:quartic}
   P(z)=z^4+az^3+bz^2+z.
\end{equation}
Then
\begin{equation}\label{eq:quartic-q}
   \widehat q(z)=\frac14\widetilde q(z),
   \qquad \widetilde q(z)=az^2+2bz+3.
\end{equation}
At critical points, this simplifies to
\begin{equation}\label{eq:quartic-simplify}
   \widetilde q(c)=2\bigl(1-c^2(a+2c)\bigr).
\end{equation}
Hence the quartic case of Smale's conjecture is equivalent to the following compact cubic statement: for every $(a,b)\in\C^2$, the cubic
\begin{equation}\label{eq:quartic-critical}
   4c^3+3ac^2+2bc+1=0
\end{equation}
has a root $c$ satisfying
\begin{equation}\label{eq:quartic-target}
   |1-c^2(a+2c)|\le\frac32.
\end{equation}

The symmetric extremal $P(z)=z^4-z$ gives equality.  The product of the three values $\widetilde q(c_k)$ is explicitly
\begin{equation}\label{eq:quartic-resultant}
   \prod_{k=1}^{3}\widetilde q(c_k)=4a^3-a^2b^2-18ab+4b^3+27.
\end{equation}
This is the most concrete part of the program: it isolates a cubic problem with a simple target disk.  A self-contained proof of \eqref{eq:quartic-target} would make the quartic slice an independent short paper.

\partline{Part VI}{The centered quintic slice}

For the centered quintic slice
\begin{equation}\label{eq:quintic}
   P(z)=z^5+az^3+bz^2+z,
\end{equation}
the critical equation has four roots $c_j(a,b)$.  The reduced quotient values satisfy a resultant product formula
\begin{equation}\label{eq:quintic-product}
   \prod_{j=1}^{4}\widehat q(c_j)=\frac{N(a,b)}{625},
\end{equation}
with
\begin{equation}\label{eq:Nab}
   N(a,b)=16a^4-4a^3b^2-128a^2+144ab^2-27b^4+256.
\end{equation}
At $(a,b)=(0,0)$, i.e. $P(z)=z^5+z$, the critical points satisfy $c_j^4=-1/5$.  Here $\widetilde q(c_j)=4$ and therefore
\begin{equation}\label{eq:quintic-extremal}
   |\widehat q(c_j)|=\frac45.
\end{equation}
This notation separates the unreduced polynomial $\widetilde q$ from the Smale-normalized quotient $\widehat q=\widetilde q/5$.

The centered quintic slice remains programmatic.  Local symbolic certificates, asymptotic descent in large parameters, and grid experiments are useful evidence, but they are not substitutes for a global proof.  The safe conclusion is that the slice is reduced to an explicit two-parameter algebraic problem.  It should be developed separately only when the local certificates or numerical data are written in reproducible detail.

\partline{Part VII}{Algorithmic branch: what is now rigorous elsewhere}

The MVC branch studies critical targets.  The algorithmic branch uses finite slopes for root finding.  The base update is \eqref{eq:pandrosion-step}; the practical difficulty is not the formula, but the choice of anchor, probe, chart, and scale.

A derivative-free descent probe in one variable is
\begin{equation}\label{eq:probe-quotient}
   Q_{\varepsilon,u}(z_0)=\frac{P(z_0+\varepsilon u)-P(z_0)}{\varepsilon u},
   \qquad u\in\{1,-1,i,-i\}.
\end{equation}
The modulus-gradient identity
\begin{equation}\label{eq:modgrad}
   \frac{d}{dt}\bigg|_{t=0}|P(z_0-tv)|^2
   =-2\operatorname{Re}\!\bigl[v\overline{P(z_0)}P'(z_0)\bigr]
\end{equation}
explains why such probes can detect descent directions.

\status{The rigorous convergence and complexity claims for finite-slope Pandrosion should be cited from the companion finite-slope theorem paper, not asserted here.  In particular, conditional local contraction, auto-probe quadratic convergence, vectorized arithmetic cost, and conditioned global basin complexity are treated there.  This part records the algorithmic intuition only.}

This resolves the main ambiguity of the earlier version.  The present manuscript no longer claims to improve known average-polynomial homotopy algorithms such as Beltran--Pardo or Lairez.  It proposes a different geometric route, to be judged by the companion theorem and by reproducible experiments.

\partline{Part VIII}{Starts, spirals, and low-discrepancy geometry}

The next question is where to start.  A rigid Cauchy ring can resonate with the root configuration.  A simple alternative is a shrinking irrational spiral:
\begin{equation}\label{eq:strategy-B}
   z_k=\rho q^k e^{ik\varphi},
   \qquad q\in(0,1),\qquad \varphi=\pi(3-\sqrt5).
\end{equation}
The radius decreases geometrically, while the golden angle distributes phases with low discrepancy.

\begin{figure}[h]
\centering
\setlength{\unitlength}{0.8mm}
\begin{picture}(78,58)
  \put(8,29){\vector(1,0){62}}
  \put(39,5){\vector(0,1){48}}
  \put(39,29){\circle{50}}
  \put(39,29){\circle{32}}
  \put(39,29){\circle{17}}
  \put(64,29){\circle*{2}}\put(20,46){\circle*{2}}\put(43,9){\circle*{2}}
  \put(52,43){\circle*{2}}\put(21,24){\circle*{2}}\put(53,23){\circle*{2}}
  \put(33,42){\circle*{2}}\put(33,16){\circle*{2}}\put(49,31){\circle*{2}}
  \put(30,30){\circle*{2}}\put(43,22){\circle*{2}}\put(41,35){\circle*{2}}
  \put(18,2){golden-angle shrinking starts}
\end{picture}
\caption{A spiral start family avoids some phase resonances of rigid rings.  This is a sampling principle, not a theorem by itself.}
\end{figure}

In computational versions of the program, this start geometry is combined with homothetic scaling, Riemann or Mobius charts, and finite-slope probes.  The empirical claim is modest: these choices enlarge the observed basin mass for many random polynomial families.  A proof requires quantitative basin estimates, such as the conditioned basin bounds in the finite-slope companion paper.

\partline{Part IX}{Extended line search and the undershooting gap}

On univariate Kostlan--Smale inputs, a Newton-like step from a typical point near unit scale has magnitude heuristically of order
\begin{equation}\label{eq:ks-newton-mag}
   |P(z)/P'(z)|\asymp d^{-1/2},
\end{equation}
whereas nearby roots are often at unit-scale distance.  Backtracking can only shrink a step that may already be too small.  This creates an undershooting gap.

The proposed repair is symmetric line search.  Instead of testing only $\tau\le1$, define
\begin{equation}\label{eq:ELS}
   \mathcal T_{\mathrm{ELS}}=\{2^k:k_-\le k\le k_+\},
\end{equation}
and choose
\begin{equation}\label{eq:ELS-step}
   z_{n+1}=z_n-\tau^\star\Delta_n,
   \qquad
   \tau^\star=\arg\min_{\tau\in\mathcal T_{\mathrm{ELS}}}|P(z_n-\tau\Delta_n)|.
\end{equation}
The predicted useful scale is $\tau\approx\sqrt d$, matching \eqref{eq:ks-newton-mag}.  This is a mechanism, not a complexity theorem.

A more careful statement is the following conjectural target.
\begin{conjecture}[Programmatic basin-mass target]
For a suitable Pandrosion start geometry on a normalized random polynomial ensemble, the probability $\beta_d$ of hitting a certified basin is bounded below by an inverse polynomial in $d$.  Under this basin-mass hypothesis, the finite-slope companion theorem yields polynomial expected sampling complexity.
\end{conjecture}

This formulation replaces the earlier unsupported $O(d\log d)$ claim.  It says exactly what remains to be proved: a lower bound on basin mass for the chosen geometric sampling distribution.

\section*{Closing view}
\addcontentsline{toc}{section}{Closing view}
The nine parts are one path through the same object.  The quotient $Q$ begins as a secant slope.  It becomes a root-finding operator, then a language for critical points, then an inverse-fiber geometry, and finally a guide for starts, scaling, and line search.

The proved content in this note is algebraic: telescoping quotients, product formulae, vanishing identities, and low-degree reductions.  The global MVC program remains open beyond the established slices and reductions.  The algorithmic branch should now be read through the finite-slope theorem and its conditioned basin-complexity companion, where the analytic claims are made precisely.

In this revised form, \emph{Universitas Pandrosion} is not offered as a proof of Smale's Mean Value Conjecture or Smale's 17th problem.  It is a compact research program: finite slopes provide the common language; fibers and critical points provide the Smale geometry; scaling and Riemann charts provide the computational geometry.

\begin{thebibliography}{15}
\bibitem{smale1981} S. Smale, \emph{The fundamental theorem of algebra and complexity theory}, Bull. Amer. Math. Soc. 4 (1981), 1--36.
\bibitem{smale1998} S. Smale, \emph{Mathematical problems for the next century}, Math. Intelligencer 20 (1998), 7--15.
\bibitem{shubsmale} M. Shub and S. Smale, \emph{Complexity of Bezout's theorem I: Geometric aspects}, J. Amer. Math. Soc. 6 (1993), 459--501.
\bibitem{beltranpardo} C. Beltran and L. M. Pardo, \emph{Smale's 17th problem: average polynomial time to compute affine and projective solutions}, J. Amer. Math. Soc. 22 (2009), 363--385.
\bibitem{lairez} P. Lairez, \emph{A deterministic algorithm to compute approximate roots of polynomial systems in polynomial average time}, Found. Comput. Math. 17 (2017), 1265--1292.
\bibitem{ortegarheinboldt} J. M. Ortega and W. C. Rheinboldt, \emph{Iterative Solution of Nonlinear Equations in Several Variables}, Academic Press, 1970; SIAM Classics edition, 2000.
\bibitem{potraptak} F.-A. Potra and V. Ptak, \emph{Nondiscrete Induction and Iterative Processes}, Pitman, 1984.
\bibitem{krawczyk} R. Krawczyk, \emph{Newton-Algorithmen zur Bestimmung von Nullstellen mit Fehlerschranken}, Computing 4 (1969), 187--201.
\bibitem{neumaier} A. Neumaier, \emph{Interval Methods for Systems of Equations}, Cambridge University Press, 1990.
\bibitem{edelmankostlan} A. Edelman and E. Kostlan, \emph{How many zeros of a random polynomial are real?}, Bull. Amer. Math. Soc. 32 (1995), 1--37.
\bibitem{crane} E. T. Crane, \emph{Extremal polynomials in Smale's mean value conjecture}, Comput. Methods Funct. Theory 6 (2006), 145--163.
\bibitem{beardonminda} A. F. Beardon, D. Minda and T. W. Ng, \emph{Smale's mean value conjecture and the hyperbolic metric}, Math. Ann. 322 (2002), 623--632.
\bibitem{bertini} D. J. Bates, J. D. Hauenstein, A. J. Sommese and C. W. Wampler, \emph{Numerically Solving Polynomial Systems with Bertini}, SIAM, 2013.
\bibitem{homcont} P. Breiding and S. Timme, \emph{HomotopyContinuation.jl: a package for homotopy continuation in Julia}, Mathematical Software -- ICMS 2018, Lecture Notes in Computer Science 10931, 458--465, Springer, 2018.
\bibitem{besevicfinite} I. Besevic, \emph{A finite-slope theorem for Pandrosion: contraction, basin mass, and conditioned complexity}, companion manuscript, 2026.
\end{thebibliography}

\end{document}
